\documentclass[10pt]{scrartcl}
\usepackage{tritricks}
%%
%
% This is a poster template with latex macros and using
% the University of Florida Logo.  For further information
% on making postscript, resizeing, and printing the poster file
% please see web site
% http://www.phys.ufl.edu/~pjh/posters/poster_howto.html
%
% N.B. This format is cribbed from one obtained from the University
% of Karlsruhe, so some macro names and parameters are in German
% Here is a short glossary:
% Breite: width
% Hoehe: height
% Spalte: column
% Kasten: box
%
% All style files necessary are part of standard TeTeX distribution
% On the UF unix cluster you should not need to import these files
% specially, as they will be automatically located.  If you
% run on a local PC however, you will need to locate these files.
% At UF try /usr/local/TeTeX...
%
% P. Hirschfeld 2/11/00
%
% The recommended procedure is to first generate a ``Special Format" size poster
% file, which is relatively easy to manipulate and view.  It can be
% resized later to A0 (900 x 1100 mm) full poster size, or A4 or Letter size
% as desired (see web site).  Note the large format printers currently
% in use at UF's OIR have max width of about 90cm or 3 ft., but the paper
% comes in rolls so the length is variable.  See below the specifications
% for width and height of various formats.  Default in the template is
% ``Special Format",  with 4 columns.
%%
%%
%% Choose your poster size:
%% For printing you will later RESIZE your poster by a factor
%%        2*sqrt(2) = 2.828    (for A0)
%%        2         = 2.00     (for A1)
%%
%%
\def\breite{450mm}     % Special Format.
\def\hoehe{319.2mm}      % Scaled by 2.82 this gives 110cm x 90cm
\def\anzspalten{4}
%%
%%\def\breite{420mm}     % A3 LANDSCAPE
%%\def\hoehe{297mm}
%%\def\anzspalten{4}
%%
%% \def\breite{297mm}     % A3 PORTRAIT
%% \def\hoehe{420mm}
%% \def\anzspalten{3}
%%
%% \def\breite{210mm}     % A4 PORTRAIT
%% \def\hoehe{297mm}
%% \def\anzspalten{2}
%%
%%
%% Procedure:
%%   Generate poster.dvi with latex
%%   Check with Ghostview
%%   Make a .ps-file with ``dvips -o poster.ps poster''
%%   Scale it with poster_resize poster.ps S
%%   where S is scale factor
%%     for Special Format->A0 S= 2.828 (= 2^(3/2)))
%%     for Special Format->A1 S= 2 (= 2^(2/2)))
%%
%% Sizes (European:)
%%   A3: 29.73 X 42.04 cm
%%   A1: 59.5 X 84.1 cm
%%   A0: 84.1 X 118.9 cm
%%   N.B. The recommended procedure is ``Special Format x 2.82"
%%   which gives 90cm x 110cm (not quite A0 dimensions).
%%
%% --------------------------------------------------------------------------
%%
%% Load the necessary packages
%%
\usepackage{palatino}
\usepackage[latin1]{inputenc}
\usepackage{epsf}
\usepackage{graphicx,psfrag,color,pstcol,pst-grad}
\usepackage{amsmath,amssymb}
\usepackage{latexsym}
\usepackage{calc}
\usepackage{multicol}
\usepackage{ifthen}
%%
%% Define the required numbers, lengths and boxes
%%
\newsavebox{\dummybox}
\newsavebox{\spalten}
%\input psfig.sty

%%
%%
\newlength{\bgwidth}\newlength{\bgheight}
\setlength\bgheight{\hoehe} \addtolength\bgheight{-1mm}
\setlength\bgwidth{\breite} \addtolength\bgwidth{-1mm}

\newlength{\kastenwidth}

%% Set paper format
\setlength\paperheight{\hoehe}
\setlength\paperwidth{\breite}
\special{papersize=\breite,\hoehe}

\topmargin -1in
\marginparsep0mm
\marginparwidth0mm
\headheight0mm
\headsep0mm


%% Minimal Margins to Make Correct Bounding Box
\setlength{\oddsidemargin}{-2.44cm}
\addtolength{\topmargin}{-3mm}
\textwidth\paperwidth
\textheight\paperheight

%%
%%
\parindent0cm
\parskip1.5ex plus0.5ex minus 0.5ex
\pagestyle{empty}




\definecolor{recoilcolor}{rgb}{1,0,0}
\definecolor{occolor}{rgb}{0,1,0}
\definecolor{pink}{rgb}{0,1,1}





\def\UberStil{\normalfont\sffamily\bfseries\large}
\def\UnterStil{\normalfont\sffamily\small}
\def\LabelStil{\normalfont\sffamily\tiny}
\def\LegStil{\normalfont\sffamily\tiny}

%%
%% Define some commands
%%
\definecolor{JG}{rgb}{0.1,0.9,0.3}

\newenvironment{kasten}{%
  \begin{lrbox}{\dummybox}%
    \begin{minipage}{0.96\linewidth}}%
    {\end{minipage}%
  \end{lrbox}%
  \raisebox{-\depth}{\psshadowbox[framesep=1em]{\usebox{\dummybox}}}\\[0.5em]}
\newenvironment{spalte}{%
  \setlength\kastenwidth{1.2\textwidth}
  \divide\kastenwidth by \anzspalten
  \begin{minipage}[t]{\kastenwidth}}{\end{minipage}\hfill}

\renewcommand{\emph}[1]{{\color{red}\textbf{#1}}}


\def\op#1{\hat{#1}}
\begin{document}
%%%%%%%%%%%%%%%%%%%%%%%%%%%%%%%%%%%%%%%%%%%%%%%%%%%%
%%%               Background                     %%%
%%%%%%%%%%%%%%%%%%%%%%%%%%%%%%%%%%%%%%%%%%%%%%%%%%%%
{\newrgbcolor{gradbegin}{0.5 0.5 1}%
  \newrgbcolor{gradend}{1 1 1}%{1 1 0.5}%
  \psframe[fillstyle=gradient,gradend=gradend,%
  gradbegin=gradbegin,gradmidpoint=0.1](\bgwidth,-\bgheight)}
\vfill
%%%%%%%%%%%%%%%%%%%%%%%%%%%%%%%%%%%%%%%%%%%%%%%%%%%%
%%%                     Header                   %%%
%%%%%%%%%%%%%%%%%%%%%%%%%%%%%%%%%%%%%%%%%%%%%%%%%%%%
\hfill
\psshadowbox{\makebox[0.95\textwidth]{%

	 \parbox[c]{0.1\linewidth}{~}
   \hfill
    \parbox[c]{0.1\linewidth}{\includegraphics[width=6cm]{left.eps}}

    \hfill
    \parbox[c]{0.70\linewidth}{%
      \begin{center}
      {
        \textbf{\Huge{MutliPipes: The Next Generation of HyperPipes}}\\[0.4em]
        \textsc{\normalsize Adam Brady, Aaron Riesbeck, Tim Menzies}\\[0.2em]
        {\normalsize Lane Department of Electrical Engineering and Computer Science, West Virginia University, Morgantown, WV 26506-6109, USA}\\[0.0em]
      }
      \end{center}}
    \hfill
    \parbox[c]{0.15\linewidth}{\includegraphics[width=2.4cm]{right.eps}}


  
\hfill}}\hfill\mbox{}\\[1.cm]
%\vspace*{1.3cm}
\begin{lrbox}{\spalten}
  \parbox[t][\textheight]{1.3\textwidth}{%
    \vspace*{0.2cm}
    \hfill
%%%%%%%%%%%%%%%%%%%%%%%%%%%%%%%%%%%%%%%%%%%%%%%%%%%%
%%%                 first column                 %%%
%%%%%%%%%%%%%%%%%%%%%%%%%%%%%%%%%%%%%%%%%%%%%%%%%%%%
    \begin{spalte}   
		\begin{kasten}
    \section*{ \hspace{0.1cm} {\color{red} \underline{SUMMARY}}}
    \Large{

\hspace{6mm}This paper explores classifiction with disjunctive sets using a modified form of HyperPipes [1] called MultiPipes. Rather than apply HyperPipes to it's intended sparse datasets, we find that it's application to non-sparse, many-class datasets typically results in several tied classification scores which we then union into a disjunction. This union presents interesting possibilities in it's high accuracy in containing the target class. Although we initially cannot predict single classes, we find that these disjunctions often eliminate large portions of possible classes. Essentialy we aren't certain what the class is, but we are very certain of what the class is not. The rest of this study explored two alternative strategies with MultiPipes. The first involved methods of reducing the disjunctive sets to single classifications. The second considered growing the disjunctive sets to optimize the accuracy of containment vs. set size.
    }
\end{kasten}

\begin{kasten}
  \section*{ \hspace{0.1cm} {\color{red} \underline{HYPERPIPES}}}
\large{
HyperPipes is a learner originally designed by Witten [3] and implemented by Eisenstein et al [1] for extremely large, sparse datasets. Rather than maintain a large working memory of statistics on each row of data, HyperPipes maintains a small data structure for each class that merely "remembers" whether a particular attribute has been encountered before. For numerics, a range of maximum and minimum values encountered is kept. When classifying, original HyperPipes classifies a row based on which class most "contains" the current attributes. For numeric attributes, a new instance is "contained" if it falls within the maximum and minimum values seen so far.

\vspace{3 mm}

Hyperpipes works well for sparse datsets, as your working memory need only contain one HyperPipe for each possible class, and each pipe must maintain only the unique symbols encountered so far along with two numeric bounds for each column. The result is a fast, dumb, scalable learner.

\vspace{3 mm}

One caveat remains in the HyperPipes algorithm. For large, sparse datasets there are enough unique columns to promote a wide variance in which HyperPipe best "contains" a class. However, for more traditional datasets with fewer columns HyperPipes' accuracy breaks down. By nature HyperPipes is strongly succeptible to outlier data, as the frequency of an encountered attribute is ignored an extremely high or low numeric will stretch the bounds for that attribute.

}
\end{kasten}

\begin{kasten}
    \section*{ \hspace{0.1cm} {\color{red} \underline{MULTIPIPES}}}
    \large{
MultiPipes was spawned from HyperPipes in an attempt to fix the issues with HyperPipes which are later described. The most notable change was the conversion of HyperPipes to a class "narrower" as opposed to a classifier. This means we anticipate that MultiPipes will return a list of potential classes rather than a single class. The pseudocode for MultiPipes is shown in the next section. In addition we were able to implement an alpha value which allows us to increase or decrease the returned set size. This allows the user to decide between high accuracy with a larger returned set size or lower accuracy with a smaller returned set size.
}
\end{kasten}
    
    \end{spalte} 
%%%%%%%%%%%%%%%%%%%%%%%%%%%%%%%%%%%%%%%%%%%%%%%%%%%%
%%%               second column                  %%%
%%%%%%%%%%%%%%%%%%%%%%%%%%%%%%%%%%%%%%%%%%%%%%%%%%%%
    \begin{spalte}   
		\begin{kasten}
  \section*{ \hspace{0.1cm} {\color{red} \underline{MULTIPIPES PSEUDOCODE}}}
  \large{
   \begin{algorithmic}
\Procedure {RunHyperPipes}{$Nothing$}
\State $HyperPipes := array[]$
\State $Guessed := 0$
\State $GuessedCorrect := 0$
\ForAll {Line in DataFile}
\State $Guessed++$
\State $GuessedCorrect \mbox{+=} Classify(MyHyperPipes,Line)$
\State $HyperPipes := AddExperience(Line,HyperPipes)$
\EndFor
\State $Accuracy := GuessedCorrect/Guessed$
\EndProcedure
\end{algorithmic}

\vspace{3 mm}

\begin{algorithmic}
\Procedure {AddExperience}{$HyperPipes$,$Line$}
\State{$Pipe := FindPipe(Line.class,HyperPipes)$}
\If{$null\; Pipe$}
\State{$HyperPipes := HyperPipes + CreateHyperPipe(Line.class)$}
\State{$Pipe := FindPipe(Line.class,HyperPipes)$}
\EndIf
\ForAll{Attr in Line}
\State{$Pipe.Attr.Num++$}
\State{$Pipe.Attr.Sum += Attr$}
\State{$Pipe.Attr.SumSquares += Attr^2$}
\State{$Pipe.Attr.Mean := Pipe.Attr.Sum/Pipe.Attr.Num$}
\State{$Temp := (Pipe.Attr.SumSquares/Pipe.Attr.Num)-Pipe.Attr.Mean^2$}
\State{$Pipe.Attr.StandardDev := sqrt(Temp)$}
\If{$Pipe.Attr.StandardDev = 0\; ||\; (abs(Pipe.Attr.Mean-Attr)/Pipe.Attr.StandardDev) < 1.96\; ||\; Pipe.Attr.Num < 10$}
\State{$Pipe.Attr.max := max(Pipe.Attr.max,Attr)$}
\State{$Pipe.Attr.min := min(Pipe.Attr.min,Attr)$}
\EndIf
\EndFor
\EndProcedure
\end{algorithmic}

\vspace{3 mm}

\begin{algorithmic}
\Procedure {Classify}{$HyperPipes$,$Line$}
\State{$BestScore := 0$}
\State{$BestClass := []$}
\ForAll{HPipe in HyperPipes}
\State{$HPipeScore := 0$}
\ForAll{Attr in Line}
\If{$Attr <= HPipe.Attr.max\; \&\&\; Attr >= HPipe.Attr.min$}
\State{$HPipeScore++$}
\EndIf
\EndFor
\If{$HPipeScore >= BestScore$}
\If{$HPipeScore := BestScore$}
\State{$BestClass := BestClass + HPipe.class$}
\Else
\State{$BestClass := array(HPipe.class)$}
\State{$BestScore := HPipeScore$}
\EndIf
\EndIf
\EndFor
\If{$BestClass\; contains\; Line.class$}
\State \Return 1
\EndIf
\State \Return 0
\EndProcedure
\end{algorithmic}
  }
\end{kasten}
 
    \end{spalte} 
%%%%%%%%%%%%%%%%%%%%%%%%%%%%%%%%%%%%%%%%%%%%%%%%%%%%
%%%               third column                   %%%
%%%%%%%%%%%%%%%%%%%%%%%%%%%%%%%%%%%%%%%%%%%%%%%%%%%%
    \begin{spalte}
		\begin{kasten}
    \section*{ \hspace{0.1cm} {\color{red} \underline{ISSUES WITH HYPERPIPES}}}
    \large{
      \section{Tied Classes}
During development of the original HyperPipes. implementation it was noted 
that the code only kept track of the most recently seen class with a score 
equal to or greater than the largest score seen to that point. This means
that when running HyperPipes if multiple classes tied in score the class
tested last would be blindly chosen. We found that this anomoly caused our 
HyperPipes implementation to prefer the last classes discovered once it 
learns many rows.
\section{Susceptible to Over Fitting}
HyperPipes has the ability to have very low overhead costs in terms of 
memory. However, this benefit is overshadowed by HyperPipes susceptibility 
to overfitting. As more and more rows are added as experience to a HyperPipe 
its min and max values with slowly approach the min and max of the attribute 
alone regardless of class. This causes HyperPipes guessing ability to 
completely fall apart. As the min and max values for the attributes in every
HyperPipe expand the scores become very high and the likelyhood of 
encountering the issue in the previous sub-section becomes greater. That is, 
as each of the HyperPipes expand their min and max attribute values start to 
become so similar to eachother that they all begin to acheive common scores.
\section{Outliers Ruin Scoring}
As shown in the AddExperience function (Program \ref{pro:AddExperience}) a HyperPipe contains a set 
of bounds for each attribute. These bounds indicate the min and a max value 
ever seen for this attribute for this class. This becomes a problem when an 
outlier for that attribute occurs. Imagine after 10,000 rows your min is 23 
and your max is 45. You encounter a new row where this attributes value is 
431. This causes your new max value to be 431. This can be a major pitfall 
if this is the only time a number this large is encountered for this class. 
 You have now expanded this HyperPipe to a max so large that this attribute 
loses its ability to classify (It always matched on this attribute for this 
class).

    }
\end{kasten}

\begin{kasten}
    \section*{ \hspace{0.1cm} {\color{red} \underline{MODEL INPUTS}}}
    \large{
    POM2 accepted eight input variables.  \textit{Criticality}: affects 
requirements development cost.  \textit{Criticality Modifier}: affects 
the severity of criticality. \textit{Culture}: affects requirements 
reordering and perceived value.  \textit{Initial-known}: sets the
size of the initial requirements set.  \textit{Inter-dependency}: 
frequency of inter-team dependency for requirements.  \textit{Size}: 
size of project (number of requirements).  \textit{Team size}: size of teams 
working on project.  \textit{Dynamism}: how frequently requirements' values 
change.

    }
\end{kasten}

\begin{kasten}
    \section*{ \hspace{0.1cm} {\color{red} \underline{RESULTS}}}
    \tiny
    
\begin{tabular}{@{ } l @{ } r | @{ } r | @{ } r | @{ } r | @{ } r | @{ } r | @{ } r | @{ } r | @{ } r | @{ } r | @{ } r |}
% \multicolumn{2}{c|}{~}&\multicolumn{10}{c}{Note: ranges run $bin_{i} <= value < bin_{i+1}$}\\
%\multicolumn{2}{c|}{~} \\
%\multicolumn{2}{c|}{~} \\
\cline{3-12}
  	&& bin 1 & bin 2 & bin 3 & bin 4 & bin 5 & bin 6 & bin 7 & bin 8 & bin 9 & bin 10 \\
          \cline{2-12}
%          & & & & & & & & & & & \\
	  & criticality
	  & .82
	  & .86
	  & .91
	  & .95
	  & 1.0
	  & 1.04
	  & 1.08
	  & 1.13
	  & 1.17
	  & 1.22\\
	  & criticality modifier
  	  & 2.0
	  & 2.8
	  & 3.6
	  & 4.4
	  & 5.2
	  & 6.0
	  & 6.8
	  & 7.6
	  & 8.4
	  & 9.2\\
	  &culture
	  & 0
	  & 10
	  & 20
	  & 30
	  & 40
	  & 50
	  & 60
	  & 70
	  & 80
	  & 90 \\
	  &initial known
	  & .40
	  & .43
	  & .46
	  & .49
	  & .52
	  & .55
	  & .58
	  & .61
	  & .64
	  & .67\\
	  &inter-dependency
	  & 0
	  & 10
	  & 20
	  & 30
	  & 40
	  & 50
	  & 60
	  & 70
	  & 80
	  & 90\\
	  &size (total personnel)
	  & 3.0
	  & 32.7
	  & 62.4
	  & 92.1
	  & 121.8
	  & 151.5
	  & 181.2
	  & 210.9
	  & 240.6
	  & 270.3\\
policy / dynamism&team size (people per team)
	  & 1.0
	  & 5.1
	  & 9.2
	  & 13.3
	  & 17.4
	  & 21.5
	  & 25.6
	  & 29.7
	  & 33.8
	  & 37.9\\
          

%policy / dynamism \\
\hline
plan-based / very low&criticality&
 &
 &
 &
 &
 &
 &
 &
 &
 &
 \\
$\sigma=\lambda=0$&criticality modifier&
 &
 &
 &
 &
 &
 &
 &
 &
 &
 \\
&culture&
 &
 &
 &
 &
 &
 &
 &
 &
 &
 \\
&initial known     &
 &
 &
 &
 &
 &
 &
 &
 &
\sq{2}{98} &
\sq{4}{96} \\
&inter-dependency     &
 &
 &
 &
 &
 &
 &
\sq{1}{99} &
\sq{1}{99} &
\sq{1}{99} &
 \\
&size     &
\sq{89}{11} &
\sq{2}{98} &
 &
 &
 &
 &
 &
 &
 &
 \\
&team size     &
 &
 &
 &
 &
 &
 &
 &
 &
 &
\sq{1}{99} \\
          \hline
 plan-based / medium&criticality     &
 &
 &
 &
 &
 &
 &
 &
\sq{5}{95} &
\sq{21}{79} &
\sq{52}{48} \\
$\sigma = 0.15, \lambda = 0.0.15$&criticality modifier&
 &
 &
 &
 &
 &
 &
 &
 &
 &
 \\
&culture&
 &
 &
 &
 &
 &
 &
 &
 &
 &
 \\
&initial known&
 &
 &
 &
 &
 &
 &
 &
 &
 &
 \\
&inter-dependency     &
 &
 &
 &
 &
 &
\sq{1}{99} &
\sq{2}{98} &
\sq{6}{94} &
\sq{6}{94} &
\sq{8}{92} \\
&size&
 &
 &
 &
 &
 &
 &
 &
 &
 &
 \\
&team size     &
 &
 &
 &
 &
 &
 &
 &
 &
 &
\sq{1}{99} \\
          \hline
plan-based / very-high &criticality     &
 &
 &
 &
 &
 &
 &
\sq{1}{99} &
\sq{6}{94} &
\sq{22}{78} &
\sq{58}{42} \\
$\sigma=2, \lambda=0.2$&criticality modifier&
 &
 &
 &
 &
 &
 &
 &
 &
 &
 \\
&culture&
 &
 &
 &
 &
 &
 &
 &
 &
 &
 \\
&initial known&
 &
 &
 &
 &
 &
 &
 &
 &
 &
 \\
&inter-dependency     &
 &
 &
 &
 &
 &
\sq{2}{98} &
\sq{3}{97} &
\sq{6}{94} &
\sq{10}{90} &
\sq{11}{89} \\
&size&
 &
 &
 &
 &
 &
 &
 &
 &
 &
 \\
&team size     &
 &
 &
 &
 &
 &
 &
 &
 &
 &
\sq{1}{99} \\
          \hline
agile 2 / very low&criticality&
 &
 &
 &
 &
 &
 &
 &
 &
 &
 \\
$\sigma=\lambda=0$&criticality modifier&
 &
 &
 &
 &
 &
 &
 &
 &
 &
 \\
&culture&
 &
 &
 &
 &
 &
 &
 &
 &
 &
 \\
&initial known     &
 &
 &
 &
 &
 &
 &
\sq{1}{99} &
\sq{4}{96} &
\sq{12}{88} &
\sq{27}{73} \\
&interdependency     &
 &
 &
 &
 &
 &
\sq{1}{99} &
\sq{4}{96} &
\sq{6}{94} &
\sq{5}{95} &
\sq{3}{97} \\
&size     &
\sq{72}{28} &
\sq{1}{99} &
 &
 &
 &
 &
 &
 &
 &
 \\
&team size     &
 &
 &
 &
 &
 &
 &
 &
 &
\sq{1}{99} &
\sq{4}{96} \\
\hline
agile 2 / medium&criticality     &
\sq{1}{99} &
 &
 &
 &
 &
 &
\sq{1}{99} &
\sq{1}{99} &
\sq{1}{99} &
\sq{1}{99} \\
$\sigma=0.15, \lambda=0.015$&criticality modifier     &
 &
 &
 &
 &
 &
 &
\sq{1}{99} &
\sq{1}{99} &
\sq{1}{99} &
\sq{1}{99} \\
&culture     &
 &
 &
 &
 &
 &
\sq{3}{97} &
\sq{10}{90} &
\sq{19}{81} &
\sq{22}{78} &
\sq{29}{71} \\
&initial known     &
 &
 &
 &
 &
\sq{1}{99} &
 &
\sq{2}{98} &
\sq{2}{98} &
\sq{2}{98} &
\sq{2}{98} \\
&interdependency     &
 &
 &
 &
\sq{1}{99} &
\sq{1}{99} &
\sq{1}{99} &
\sq{1}{99} &
 &
\sq{1}{99} &
 \\
&size     &
\sq{97}{3} &
 &
 &
 &
 &
 &
 &
 &
 &
 \\
&team size     &
 &
\sq{17}{83} &
\sq{20}{80} &
\sq{11}{89} &
\sq{6}{94} &
\sq{1}{99} &
 &
 &
 &
 \\
          \hline
agile 2 / very high&criticality     &
\sq{1}{99} &
 &
 &
 &
\sq{1}{99} &
 &
\sq{1}{99} &
\sq{1}{99} &
\sq{1}{99} &
\sq{1}{99} \\
$\sigma=2, \lambda=0.2$&criticality modifier     &
\sq{1}{99} &
\sq{1}{99} &
\sq{1}{99} &
\sq{1}{99} &
\sq{1}{99} &
\sq{1}{99} &
\sq{1}{99} &
 &
\sq{1}{99} &
\sq{1}{99} \\
&culture     &
 &
 &
 &
 &
\sq{1}{99} &
\sq{4}{96} &
\sq{13}{87} &
\sq{17}{83} &
\sq{20}{80} &
\sq{22}{78} \\
&initial known     &
 &
 &
 &
 &
 &
\sq{1}{99} &
\sq{1}{99} &
\sq{1}{99} &
\sq{1}{99} &
\sq{1}{99} \\
&interdependency     &
 &
\sq{1}{99} &
\sq{1}{99} &
\sq{1}{99} &
\sq{1}{99} &
\sq{1}{99} &
\sq{1}{99} &
\sq{1}{99} &
 &
\sq{1}{99} \\
&size     &
\sq{100}{0} &
 &
 &
 &
 &
 &
 &
 &
 &
 \\
&team size     &
 &
\sq{15}{85} &
\sq{18}{82} &
\sq{11}{89} &
\sq{5}{95} &
\sq{2}{98} &
\sq{1}{99} &
 &
 &
\\
\end{tabular}


    \vspace{-.5em}
\end{kasten}

    \end{spalte}
%%%%%%%%%%%%%%%%%%%%%%%%%%%%%%%%%%%%%%%%%%%%%%%%%%%%
%%%               fourth column                  %%%
%%%%%%%%%%%%%%%%%%%%%%%%%%%%%%%%%%%%%%%%%%%%%%%%%%%%
    \begin{spalte}
		\begin{kasten}
    \section*{ \hspace{0.1cm} {\color{red} \underline{RESULTS}}}
    \large{
To directly compare single classification and multiple classification algorithms we developed an appropriate scoring mechanism that penalizes for larger return sets. The algorithm used for the boxplots is:

\begin{equation}
Score = 1 - \frac{(ClassesReturned - 1)}{(TotalClasses - 1)}
\end{equation}



\begin{figure}[H]
\begin{center}
\renewcommand{\baselinestretch}{0.5}
\noindent
{\scriptsize
\begin{tabular}{c r  @{} c c }
\multicolumn{4}{c}{Weighted Performance (Multiple Datasets)} \\\hline

Rank & Treatment  & 50\% & \% Ret \\
\hline

1 & NaiveBayes & \boxplot{68.8}{18.6}{87.4}{6.7}{94.1} & 14.1 \\
2 & WeightedMeans1 & \boxplot{53.9}{18.7}{72.6}{12}{84.6} & 14.1 \\
3 & WeightedMeans0 & \boxplot{53.9}{18.7}{72.6}{12.6}{85.2} & 14.1 \\
4 & Centroid & \boxplot{48.1}{15.5}{63.6}{19.5}{83.1} & 14.1 \\
5 & HyperPipes & \boxplot{40.4}{14.5}{54.9}{22}{76.9} & 14.1 \\
6 & RemoveOutliers & \boxplot{42.4}{11}{53.4}{26.5}{79.9} & 60.2 \\
7 & OverfitRevert & \boxplot{38.1}{10.2}{48.3}{21.5}{69.8} & 66.4 \\
8 & MultiPipes & \boxplot{38.1}{9.2}{47.3}{22.7}{70.0} & 67.3 \\
9 & 5\%Alpha & \boxplot{2.4}{17.3}{19.7}{32.2}{51.9} & 88.4 \\
10 & 25\%Alpha & \boxplot{0.7}{2.8}{3.5}{29.8}{33.3} & 97.1 \\
11 & 50\%Alpha & \boxplot{0.2}{0.5}{0.7}{14.1}{14.8} & 99.3 \\



\end{tabular}
}
\end{center}
\end{figure}


\begin{figure}[H]
\begin{center}
\renewcommand{\baselinestretch}{0.5}
\noindent
{\scriptsize
\begin{tabular}{c r  @{} c }
\multicolumn{3}{c}{Weighted Performance (Many Classes Datasets)} \\\hline

Rank & Treatment  & 50\% \\
\hline

1 & WeightedMeans1 & \boxplot{53.7}{16.3}{70.0}{18.9}{88.9} \\
2 & WeightedMeans0 & \boxplot{53.7}{16.3}{70.0}{18.9}{88.9}\\
3 & MultiPipes & \boxplot{53.7}{16.3}{70.0}{18.9}{88.9}\\
4 & RemoveOutliers & \boxplot{53.7}{16.1}{69.8}{19.1}{88.9} \\
5 & OverfitRevert & \boxplot{53.7}{16.1}{69.8}{19.1}{88.9} \\
6 & NaiveBayes & \boxplot{42.1}{23.2}{65.3}{18.4}{83.7} \\
7 & HyperPipes & \boxplot{21.1}{38}{59.1}{16.2}{75.3} \\
8 & Centroid & \boxplot{23.4}{34.4}{57.8}{16.4}{74.2} \\
9 & 5\%Alpha & \boxplot{33.8}{20.6}{54.4}{14.7}{69.1} \\
10 & 25\%Alpha & \boxplot{22.7}{11.1}{33.8}{24.1}{57.9} \\
11 & 50\%Alpha & \boxplot{18.2}{4.4}{22.6}{3.2}{25.8} \\



\end{tabular}
}
\end{center}
\end{figure}



    }

\end{kasten}

\begin{kasten}
    \section*{ \hspace{0.1cm} {\color{red} \underline{CONCLUSIONS}}}
    \large{
Lamentably, on many-class (approx. 20) datasets Multipipes does not perform significantly better than Naive Bayes when scored using our performance measure to penalize overclassification. Even worse, when tested on datasets with between two and twenty-four classes Multipipes does significantly worse.

\vspace{3 mm}

Also, Multipipes tends to over-classify despite our efforts to reign in overfitting. On average Multipipes returns two-thirds of the possible classes, the average number of classes for a dataset being seven. While there are potential benefits from reducing the space of possible classes, the accuracy of our disjunctions is only slightly better than Naive Bayes'. Increasing the "reach" of our disjunctions using alpha values improves this accuracy to over 99\%, but at that point Multipipes is typically returning every single possible class.

\vspace{3 mm}

In conclusion, while there are potential uses in the case of very-many class systems, there appear to exist intrinsic drawbacks to Multipipes' abilities as a classifier.
    }
\end{kasten}

\begin{kasten}
    \section*{ \hspace{0.1cm} {\color{red} \underline{FOR MORE INFORMATION}}}
    \vspace{-0.5em}
    \normalsize{
      Aaron Riesbeck (\url{aaron.riesbeck@gmail.com})\\
      Adam Brady (\url{adam.m.brady@gmail.com})\\
      Tim Menzies Ph.D. (\url{tim@menzies.us})\\
    }
    \vspace{-0.5em}
\end{kasten}

\begin{kasten}
    \section*{ \hspace{0.1cm} {\color{red} \underline{References}}}
    \vspace{-0.5em}
    \normalsize{
      \begin{enumerate}
      \item J. Eisenstein and R. Davis. Visual and linguistic information in gesture classification. In ICMI, pages 113-120, 2004.
      \item R. J. Larsen and M. L. Marx. An Introduction to Mathematical Statistics and Its Applications, Third Edition. Prentice Hall, 2001.
      \item I. H. Witten and E. Frank. Data Mining: Practical Machine Learning Tools and Techniques with Java Implementations. Morgan Kaufmann, 1999.
      \end{enumerate}
    }
\end{kasten}

    \end{spalte}
    }
    \end{lrbox}
\resizebox*{0.98\textwidth}{!}{%
  \usebox{\spalten}}\hfill\mbox{}\vfill
\end{document}


